% !TEX program = xelatex
\documentclass[11pt, a4paper]{article}
\usepackage[a4paper, margin=2.5cm]{geometry}
\usepackage{fontspec}
\usepackage[portuguese, bidi=basic, provide=*]{babel}
\babelfont{rm}{Noto Sans}
\usepackage{amsmath, amssymb, amsthm}
\usepackage{tikz}
\usepackage{booktabs}
\usepackage{microtype}

% \usepackage[utf8]{inputenc}
\usepackage{enumitem}

\usetikzlibrary{shapes.geometric, arrows.meta, positioning, fit, backgrounds, decorations.pathreplacing}

\setlist[itemize]{label={--}}

\title{\textbf{Roteiro para Lista 1}}
\date{\today}

\begin{document}
\maketitle

% EXERCÍCIO 1
\subsection*{Exercício 1: Teorema de Ramsey Multicor ($K_N \to K_k$)}
\textbf{Intuição Estrutural:} A prova baseia-se na aplicação iterativa do Princípio da Casa dos Pombos para construir uma sequência de vizinhanças aninhadas, forçando uma cor dominante a cada passo.

\textbf{Passo a Passo Rigoroso:}
\begin{enumerate}
    \item \textbf{Construção Iterativa:} Dado um vértice $v_1$ em $A_0 = V(K_n)$, ele tem $n-1$ arestas. Pelo Princípio da Casa dos Pombos, pelo menos $(n-1)/r$ arestas compartilham uma mesma cor $c_1$. Define-se $A_1$ como os vizinhos alcançados por $c_1$.
    \item \textbf{Aninhamento:} Repete-se o processo tomando $v_2 \in A_1$, garantindo uma cor $c_2$ dominante para $A_2$, resultando em $A_0 \supseteq A_1 \supseteq \dots \supseteq A_t$.
    \item \textbf{Decisão Unilateral:} Para qualquer par $v_i, v_j$ (com $i < j$), a aresta $v_iv_j$ tem cor $c_i$, pois $v_j \in A_{j-1} \subseteq A_i$. A cor depende exclusivamente do vértice de menor índice.
    \item \textbf{Garantia da Clique:} Se $t = r(k-1)+1$, uma mesma cor $c$ deve aparecer entre os $c_i$ pelo menos $k$ vezes (Casa dos Pombos novamente). Os vértices correspondentes a esses passos formarão um $K_k$ monocromático.
\end{enumerate}
\begin{tikzpicture}[
    >=Stealth,
    vertex/.style={circle, draw=black, thick, fill=white, minimum size=18pt, inner sep=1pt, font=\small},
    red_edge/.style={draw=red!80!black, very thick},
    blue_edge/.style={draw=blue!80!black, very thick, dashed},
    set_box/.style={ellipse, draw=gray!70, thick, fill=gray!5}
]

    % Título da Parte 1
    \node[font=\bfseries\large] at (5.5, 4.2) {1. Construção dos Conjuntos Encaixados ($A_0 \supseteq A_1 \supseteq A_2 \supseteq A_3$)};

    % Conjuntos Encaixados (Elipses da maior para a menor)
    \draw[fill=gray!8, draw=gray!60, thick] (5.5,0) ellipse (6.2cm and 3.2cm);
    \node[font=\bfseries] at (0.3, 2.2) {$A_0 = V(K_n)$};

    \draw[fill=red!6, draw=red!50!black, thick] (6.5,0) ellipse (4.8cm and 2.5cm);
    \node[font=\bfseries, red!60!black] at (2.8, 1.8) {$A_1$};

    \draw[fill=blue!6, draw=blue!50!black, thick] (7.6,0) ellipse (3.4cm and 1.8cm);
    \node[font=\bfseries, blue!60!black] at (5.0, 1.3) {$A_2$};

    \draw[fill=red!12, draw=red!60!black, thick] (8.7,0) ellipse (2.0cm and 1.1cm);
    \node[font=\bfseries, red!60!black] at (7.4, 0.7) {$A_3$};

    % Vértices escolhidos em cada etapa
    \node[vertex, fill=red!20] (v1) at (0.5, 0) {$v_1$};
    \node[vertex, fill=blue!20] (v2) at (3.2, 0) {$v_2$};
    \node[vertex, fill=red!20] (v3) at (5.8, 0) {$v_3$};
    \node[vertex, fill=red!20] (v4) at (8.8, 0) {$v_j$};

    % Arestas de v_1 (cor c_1 = vermelho) para todo A_1
    \draw[red_edge, ->] (v1) to[bend left=25] node[above, sloped, font=\scriptsize] {$\chi(v_1 u) = c_1$} (v2);
    \draw[red_edge, ->] (v1) to[bend left=38] (v3);
    \draw[red_edge, ->] (v1) to[bend left=45] (v4);

    % Arestas de v_2 (cor c_2 = azul) para todo A_2
    \draw[blue_edge, ->] (v2) to[bend right=25] node[below, sloped, font=\scriptsize] {$\chi(v_2 u) = c_2$} (v3);
    \draw[blue_edge, ->] (v2) to[bend right=38] (v4);

    % Arestas de v_3 (cor c_3 = vermelho) para todo A_3
    \draw[red_edge, ->] (v3) to node[above, font=\scriptsize] {$\chi(v_3 u) = c_3$} (v4);

    % Legenda / Explicação na base
    \node[align=center, font=\footnotesize, draw=black!40, rounded corners, fill=yellow!10, inner sep=6pt] at (5.5, -4.2) {
        \textbf{Propriedade Fundamental:} Para $i < j$, como $v_j \in A_{j-1} \subseteq A_i$, a cor de $v_i v_j$ é decidida unicamente por $v_i$ ($\chi(v_i v_j) = c_i$).\\
        Acima, $c_1 = c_3 = \text{vermelho}$. Logo, os vértices $\{v_1, v_3, v_j\}$ formam um triângulo ($K_3$) inteiramente vermelho!
    };

\end{tikzpicture}

% EXERCÍCIO 2
\subsection*{Exercício 2: Limite Inferior de Erdős para R(k) (1947)}
\textbf{Intuição Probabilística:} Ao invés de construir uma coloração, contamos todas as colorações e mostramos que a proporção daquelas que contêm uma clique monocromática é estritamente menor que 1. Se a probabilidade de falha é $< 1$, uma coloração válida \textit{deve} existir.

\begin{tikzpicture}[
    >=Stealth,
    vertexS/.style={circle, draw=red!80!black, very thick, fill=red!15, minimum size=20pt, font=\small\bfseries},
    vertexR/.style={circle, draw=gray!80!black, thick, fill=gray!15, minimum size=20pt, font=\small},
    locked_edge/.style={draw=red!80!black, line width=1.8pt},
    free_edge/.style={draw=gray!70!black, thick, dashed},
    boxS/.style={draw=red!80!black, thick, fill=red!10, minimum width=0.8cm, minimum height=0.7cm, font=\scriptsize},
    boxF/.style={draw=gray!70!black, thick, fill=gray!10, minimum width=0.8cm, minimum height=0.7cm, font=\scriptsize}
]

    % =========================================================
    % PARTE ESQUERDA: ESTRUTURA DO GRAFO K_n
    % =========================================================
    \node[font=\bfseries\normalsize] at (0, 4.3) {1. Partição das Arestas em $K_n$};

    % Elipse do grafo K_n
    \draw[fill=gray!5, draw=gray!50, thick, rounded corners=20pt] (-3.8, -2.5) rectangle (3.8, 3.5);
    \node[font=\bfseries, gray!60!black] at (3.1, 3.1) {$K_n$};

    % Conjunto S (k vértices)
    \draw[fill=red!6, draw=red!60!black, thick, dashed] (-1.5, 0.6) circle (1.8cm);
    \node[font=\bfseries, red!75!black] at (-1.5, 2.75) {Subconjunto $S$ ($|S|=k$)};

    % Vértices de S
    \node[vertexS] (v1) at (-2.4, 1.3) {$v_1$};
    \node[vertexS] (v2) at (-0.6, 1.3) {$v_2$};
    \node[vertexS] (vk) at (-1.5, -0.4) {$v_k$};

    % Vértices fora de S (V \ S)
    \node[font=\bfseries, gray!70!black] at (2.1, 2.3) {$V(K_n) \setminus S$};
    \node[vertexR] (u1) at (2.1, 1.3) {$u_1$};
    \node[vertexR] (u2) at (2.6, -0.2) {$u_2$};
    \node[vertexR] (un) at (1.5, -1.5) {$u_{n-k}$};

    % Arestas Travadas (Internas a S)
    \draw[locked_edge] (v1) -- (v2);
    \draw[locked_edge] (v2) -- (vk);
    \draw[locked_edge] (v1) -- (vk);

    % Arestas Livres (Cruzando de S para V\S e internas a V\S)
    \draw[free_edge] (v2) -- (u1);
    \draw[free_edge] (v2) -- (u2);
    \draw[free_edge] (vk) -- (u1);
    \draw[free_edge] (vk) -- (u2);
    \draw[free_edge] (vk) -- (un);
    \draw[free_edge] (u1) -- (u2);
    \draw[free_edge] (u2) -- (un);
    \draw[free_edge] (u1) to[bend left=20] (un);

    % Legenda do Grafo
    \draw[locked_edge] (-3.2, -1.8) -- (-2.8, -1.8) node[right, font=\footnotesize, text=black] {Arestas fixadas em $S$: $\binom{k}{2}$};
    \draw[free_edge]   (-3.2, -2.2) -- (-2.8, -2.2) node[right, font=\footnotesize, text=black] {Arestas livres fora de $S$: $\binom{n}{2}-\binom{k}{2}$};

    % =========================================================
    % PARTE DIREITA: CONTAGEM MULTIPLICATIVA (VETOR DE ARESTAS)
    % =========================================================
    \node[font=\bfseries\normalsize] at (9.2, 4.3) {2. Contagem de Colorações com $S$ Monocromático};

    % Bloco Travado (internas a S)
    \node[boxS] (e1) at (5.8, 1.8) {$e_1$};
    \node[boxS, right=0pt of e1] (e2) {$e_2$};
    \node[boxS, right=0pt of e2] (edots1) {$\cdots$};
    \node[boxS, right=0pt of edots1] (ek) {$e_{\binom{k}{2}}$};

    % Bloco Livre (fora de S)
    \node[boxF, right=0.4cm of ek] (f1) {$e_{\binom{k}{2}+1}$};
    \node[boxF, right=0pt of f1] (f2) {$e_{\binom{k}{2}+2}$};
    \node[boxF, right=0pt of f2] (fdots) {$\cdots$};
    \node[boxF, right=0pt of fdots] (fn) {$e_{\binom{n}{2}}$};

    % Chaves superiores
    \draw[decorate, decoration={brace, amplitude=6pt}, thick, red!80!black]
        (e1.north west) -- (ek.north east)
        node[midway, above=8pt, font=\footnotesize\bfseries, align=center]
        {$\binom{k}{2}$ arestas internas a $S$\\(Bloco Travado)};

    \draw[decorate, decoration={brace, amplitude=6pt}, thick, gray!70!black]
        (f1.north west) -- (fn.north east)
        node[midway, above=8pt, font=\footnotesize\bfseries, align=center]
        {$\binom{n}{2} - \binom{k}{2}$ arestas restantes\\(Arestas Livres)};

    % Chaves inferiores explicando as escolhas
    \draw[decorate, decoration={brace, amplitude=6pt, mirror}, thick, red!80!black]
        (e1.south west) -- (ek.south east)
        node[midway, below=8pt, font=\tiny, align=center] (c1)
        {Todas da mesma cor:\\ \textbf{2 opções ao todo}\\(Tudo Vermelho ou Tudo Azul)};

    \draw[decorate, decoration={brace, amplitude=6pt, mirror}, thick, gray!70!black]
        (f1.south west) -- (fn.south east)
        node[midway, below=8pt, font=\tiny, align=center] (c2)
        {2 cores independentes por aresta:\\ $\underbrace{2 \times 2 \times \cdots \times 2}_{\binom{n}{2}-\binom{k}{2}\text{ vezes}} = \mathbf{2^{\binom{n}{2}-\binom{k}{2}}}$ \textbf{opções}};

    % Caixa de resultado final
    \node[draw=black!70, thick, rounded corners=6pt, fill=yellow!12, inner sep=8pt, font=\small, align=center]
        at (9.2, -1.5) {
        \textbf{Pelo Princípio Multiplicativo (para $S$ fixo):}\\[4pt]
        $\displaystyle \underbrace{2}_{\text{cor de } S} \cdot \underbrace{2^{\binom{n}{2}-\binom{k}{2}}}_{\text{arestas livres}} = 2^{\binom{n}{2}-\binom{k}{2}+1} \text{ colorações}$
    };

\end{tikzpicture}

\begin{enumerate}
    \item \textbf{Etapa 1 --- Definição do Espaço Amostral e Estrutura das Colorações:}
    \begin{itemize}
        \item Considere o grafo completo $K_n$, o qual possui exatamente $\binom{n}{2}$ arestas.
        \item Atribuímos independentemente uma de duas cores (vermelho ou azul) a cada aresta, resultando em um total de $2^{\binom{n}{2}}$ colorações distintas para $E(K_n)$.
        \item Chamamos uma coloração de \textbf{ruim} se ela contiver ao menos uma $k$-clique ($K_k$) monocromática, e de \textbf{boa} caso contrário.
    \end{itemize}

    \item \textbf{Etapa 2 --- Contagem das Colorações Ruins via Cota da União:}
    \begin{itemize}
        \item Fixe um subconjunto específico $S \subseteq V(K_n)$ composto por $|S| = k$ vértices.
        \item Para que $S$ seja monocromático: há $2$ escolhas para a cor de $S$ (tudo vermelho ou tudo azul), o que fixa as $\binom{k}{2}$ arestas internas a $S$, enquanto as demais $\binom{n}{2} - \binom{k}{2}$ arestas de $K_n$ permanecem livres, podendo ser coloridas de $2^{\binom{n}{2} - \binom{k}{2}}$ maneiras.
        \item Assim, o número de colorações que tornam $S$ monocromático é $2 \cdot 2^{\binom{n}{2} - \binom{k}{2}} = 2^{\binom{n}{2} - \binom{k}{2} + 1}$.
        \item Como existem $\binom{n}{k}$ subconjuntos de $k$ vértices em $K_n$, aplicamos o princípio multiplicativo:
        \[ \#\{\text{ruins}\} \le \binom{n}{k} 2^{\binom{n}{2} - \binom{k}{2} + 1}. \]
    \end{itemize}

    \item \textbf{Etapa 3 --- Condição para Existência de uma Coloração Boa:}
    \begin{itemize}
        \item Se o número de colorações ruins for estritamente menor que o total $2^{\binom{n}{2}}$, garante-se que pelo menos uma coloração boa sobra:
        \[ \#\{\text{ruins}\} \le \binom{n}{k} 2^{\binom{n}{2} - \binom{k}{2} + 1} < 2^{\binom{n}{2}}. \]
        \item Dividindo ambos os lados por $2^{\binom{n}{2}}$, obtemos a desigualdade equivalente:
        \[ \binom{n}{k} 2^{1 - \binom{k}{2}} < 1. \]
    \end{itemize}

    \item \textbf{Etapa 4 --- Aplicação da Cota de Stirling e Substituição de $n \le 2^{k/2}$:}
    \begin{itemize}
        \item Usando $\binom{n}{k} \le \frac{n^k}{k!}$ junto com $k! \ge (k/e)^k$, estimamos $\binom{n}{k} \le \left(\frac{en}{k}\right)^k$.
        \item Supondo $n \le 2^{k/2}$, temos $n^k \le 2^{k^2/2}$ e, portanto, $\binom{n}{k} \le \left(\frac{e}{k}\right)^k 2^{k^2/2}$.
        \item Multiplicando por $2^{1 - \binom{k}{2}} = 2^{1 - \frac{k^2}{2} + \frac{k}{2}}$ e somando os expoentes de base $2$ (onde $\frac{k^2}{2}$ se cancela), obtemos:
        \[ \binom{n}{k} 2^{1 - \binom{k}{2}} \le \left(\frac{e}{k}\right)^k 2^{\frac{k^2}{2}} \cdot 2^{1 - \frac{k^2}{2} + \frac{k}{2}} = \left(\frac{e}{k}\right)^k 2^{1 + k/2} = 2\left(\frac{e\sqrt{2}}{k}\right)^k. \]
    \end{itemize}

    \item \textbf{Etapa 5 --- Análise Numérica e Limiar $k \ge 5$:}
    \begin{itemize}
        \item Como $e\sqrt{2} \approx 3{,}844 < k$ para todo $k \ge 4$, a base $\frac{e\sqrt{2}}{k}$ é menor que $1$ e a expressão $g(k) = 2\left(\frac{e\sqrt{2}}{k}\right)^k$ é estritamente decrescente em $k$ para $k \ge 4$.
        \item Para $k = 4$, a expressão vale $\approx 1{,}706 > 1$ (a cota falha). Porém, para $k = 5$, temos:
        \[ 2\left(\frac{e\sqrt{2}}{5}\right)^5 \approx 2 \cdot (0{,}7688)^5 \approx 0{,}537 < 1. \]
        \item Pela monotonicidade decrescente, conclui-se que $2\left(\frac{e\sqrt{2}}{k}\right)^k < 1$ para todo $k \ge 5$.
    \end{itemize}

    \item \textbf{Etapa 6 --- Escolha de $n = \lfloor 2^{k/2} \rfloor$ e Conclusão sobre $R(k)$:}
    \begin{itemize}
        \item \textbf{Garantia de $n \le 2^{k/2}$ e $n \ge k$:} Como o número de vértices deve ser inteiro, definimos $n := \lfloor 2^{k/2} \rfloor$. Pela definição da função piso ($\lfloor x \rfloor \le x < \lfloor x \rfloor + 1$), temos imediatamente $n \le 2^{k/2}$. Além disso, para $k = 5$, temos $2^{5/2} = 4\sqrt{2} \approx 5{,}656 \implies n = \lfloor 5{,}656 \rfloor = 5 \ge 5$; como $2^{k/2}$ cresce exponencialmente, segue que $n = \lfloor 2^{k/2} \rfloor \ge k$ para todo $k \ge 5$, garantindo que $K_n$ possui ao menos $k$ vértices.
        \item \textbf{Existência da Coloração Boa em $K_n$:} Como $n = \lfloor 2^{k/2} \rfloor$ satisfaz $n \le 2^{k/2}$, as Etapas 4 e 5 asseguram que $\binom{n}{k} 2^{1 - \binom{k}{2}} < 1$ para todo $k \ge 5$. Multiplicando por $2^{\binom{n}{2}}$, temos $\#\{\text{ruins}\} < 2^{\binom{n}{2}}$, donde $\#\{\text{boas}\} = 2^{\binom{n}{2}} - \#\{\text{ruins}\} \ge 1$. Logo, existe pelo menos uma $2$-coloração das arestas de $K_n$ sem nenhuma clique monocromática de tamanho $k$ ($K_n \not\to K_k$), e restringir essa coloração a qualquer $m \le n$ também evita $K_k$ monocromático.
        \item \textbf{Cota Final para $R(k)$:} Sendo $R(k)$ o menor inteiro $N$ tal que toda $2$-coloração de $K_N$ contém um $K_k$ monocromático, o fato de $K_n$ admitir uma coloração boa implica que $R(k) > n = \lfloor 2^{k/2} \rfloor$. Como $R(k)$ é um número inteiro, ele deve ser pelo menos $n + 1$:
        \[ R(k) \ge \lfloor 2^{k/2} \rfloor + 1 > 2^{k/2} \quad \text{para todo } k \ge 5. \]
    \end{itemize}
\end{enumerate}

% EXERCÍCIO 3
\subsection*{Exercício 3: Teorema de Kővári-Sós-Turán}
\textbf{Intuição Estrutural:} Dupla contagem em grafos estrelas $K_{1,s}$ (chamados de $s$-cerejas). Contamos as cerejas a partir do vértice central e, depois, limitamos a partir das pontas (conjunto de $s$ vértices).

\textbf{Passo a Passo Rigoroso:}
\begin{enumerate}
    \item \textbf{Contagem a partir do centro:} Um vértice $v$ forma $\binom{d(v)}{s}$ cerejas. O total no grafo é $\sum \binom{d(v)}{s}$.
    \item \textbf{Contagem pelas extremidades:} Qualquer subconjunto de $s$ vértices só pode ser o lado direito de, no máximo, $t-1$ cerejas (se fossem $t$, formariam um $K_{s,t}$). O total é $\le (t-1)\binom{n}{s}$.
    \item \textbf{Convexidade:} Como $\binom{x}{s}$ é convexa, o mínimo ocorre no grau médio $d_{avg} = 2|E(G)|/n$.
    \[ n\binom{2|E(G)|/n}{s} \le (t-1)\binom{n}{s} \]
    \item \textbf{Simplificação Assintótica:} Expandindo os binomiais pelo termo dominante $\mathcal{O}(x^s)$, isola-se $|E(G)|$, provando que o crescimento é limitado por $\mathcal{O}(n^{2-1/s})$.
\end{enumerate}

% EXERCÍCIO 4
\subsection*{Exercício 4: Lema de Pósa (Caminhos e Ciclos)}
\textbf{Intuição Estrutural:} Em um caminho maximal, todos os vizinhos das extremidades devem estar contidos no próprio caminho. Pelo princípio da casa dos pombos, a intersecção de vizinhanças força uma aresta de retorno que cria um ciclo longo, impossibilitando a maximalidade se o caminho for curto.

\begin{center}
\begin{tikzpicture}[scale=1, every node/.style={circle, draw, minimum size=8mm, inner sep=1pt}]
    % Vértices do caminho
    \node (v0) at (0,0) {$v_0$};
    \node (v1) at (1.5,0) {$v_1$};
    \node[draw=none] (dots1) at (3,0) {$\dots$};
    \node (vi1) at (4.5,0) {$v_{i-1}$};
    \node (vi) at (6,0) {$v_i$};
    \node[draw=none] (dots2) at (7.5,0) {$\dots$};
    \node (vk) at (9,0) {$v_k$};

    % Caminho base
    \draw[thick] (v0) -- (v1) -- (dots1) -- (vi1) -- (vi) -- (dots2) -- (vk);

    % Arestas que fecham o ciclo
    \draw[thick, blue, bend right=45] (vk) to node[above, draw=none] {vizinho de $v_k$} (vi1);
    \draw[thick, red, bend left=45] (v0) to node[above, draw=none] {vizinho de $v_0$} (vi);
\end{tikzpicture}
\end{center}

\textbf{Passo a Passo Rigoroso:}
\begin{enumerate}
    \item Assuma um caminho maximal $P$ com $k < 2\delta(G)$. Os conjuntos de vizinhos $N(v_0)$ e $N(v_k)$ estão obrigatoriamente dentro de $P$.
    \item Desloque os vizinhos de $v_0$ uma posição para trás no caminho: $S = \{v_{i-1} \mid v_i \in N(v_0)\}$. O tamanho é $|S| \ge \delta(G)$.
    \item Como $|S| + |N(v_k)| \ge 2\delta(G) > k$, os conjuntos se interceptam. Existe $v_{i-1}$ que é vizinho de $v_k$, cujo sucessor $v_i$ é vizinho de $v_0$.
    \item O ciclo $C = (v_0, v_i, v_{i+1}, \dots, v_k, v_{i-1}, v_{i-2}, \dots, v_0)$ consome todos os vértices de $P$. Se o grafo é conexo e $k < n-1$, existe um vértice fora do ciclo conectado a ele, permitindo estender $P$, gerando um absurdo.
\end{enumerate}

% EXERCÍCIO 5
\subsection*{Exercício 5: Subgrafo Bipartido em Grafo Livre de Triângulos}
\textbf{Intuição Estrutural:} Se não há triângulos, a vizinhança do vértice de grau máximo $\Delta$ é perfeitamente independente. Dividimos o grafo em vizinhança e não-vizinhança.

\textbf{Passo a Passo Rigoroso:}
\begin{enumerate}
    \item Seja $v_{\Delta}$ o vértice de grau máximo. O conjunto $A = N(v_{\Delta})$ não possui arestas, senão formaria um triângulo com $v_{\Delta}$.
    \item O conjunto restante é $B = V(G) \setminus A$. O número de arestas destruidoras da bipartição está inteiramente dentro de $B$.
    \item Contagem de graus: A soma dos graus em $B$ contabiliza as arestas de $e(B)$ duas vezes e as de cruzamento uma vez. $k = |e(B)| + |e(G)| \le |V(B)|\Delta = (n-\Delta)\Delta$.
    \item Por médias (MA-MG), $\Delta(n-\Delta) \le n^2/4$. Portanto $|e(G)| + |e(B)| \le n^2/4$. Como $|e(G)| \ge n^2/4 - t$, deduz-se que $|e(B)| \le t$.
\end{enumerate}

% EXERCÍCIO 6
\subsection*{Exercício 6: Triângulos em Grafos $t$-longe de Bipartidos}
\textbf{Intuição Estrutural:} Uma aresta interna à vizinhança $B_u = N(u)$ de um vértice $u$ forma exatamente um triângulo com $u$. Ao somarmos $e(B_u)$ para todos os vértices do grafo, obtemos $3T(G)$. A prova divide-se em duas frentes: (I) uma análise local em cada vértice $u$, onde particionamos $V(G)$ em $B_u = N(u)$ e $A_u = V(G) \setminus N(u)$, usamos o fato de $G$ ser $t$-longe de ser bipartido para limitar as arestas cruzadas $e(A_u, B_u)$ e calculamos a soma dos graus em $A_u$ para isolar $e(B_u)$; e (II) uma agregação global sobre todos os vértices $u \in V(G)$, aplicando dupla contagem (inversão de somatório) e a desigualdade das médias (MA-MG) para limitar o termo quadrático $\sum d(v)(n-d(v))$.

\begin{center}
\begin{tikzpicture}[scale=0.95, every node/.style={circle, draw, minimum size=6.5mm, inner sep=1pt}]
    % Vértice raiz u (destacado no topo para visualizar o triângulo)
    \node[draw=black, fill=gray!15, very thick] (u) at (3.0, 3.6) {$u$};

    % Caixa do conjunto B_u = N(u)
    \draw[rounded corners=8pt, draw=red!75!black, thick, fill=red!4] (0, 1.4) rectangle (6.0, 2.6);
    \node[draw=none, fill=none, anchor=west] at (6.2, 2.0) {\small $B_u = N(u)$ ($|B_u| = d(u)$)};

    % Vértices em B_u
    \node[draw=red!80!black, thick, fill=white] (w1) at (1.5, 2.0) {$w_1$};
    \node[draw=red!80!black, thick, fill=white] (w2) at (3.2, 2.0) {$w_2$};
    \node[draw=red!80!black, thick, fill=white] (w3) at (4.8, 2.0) {$w_3$};

    % Triângulo formado por u e uma aresta interna de B_u
    \draw[very thick, red!80!black] (w1) -- node[above, draw=none, fill=none, font=\scriptsize] {$e(B_u)$} (w2);
    \draw[thick, red!80!black, dashed] (u) -- (w1);
    \draw[thick, red!80!black, dashed] (u) -- (w2);
    \draw[thick, gray!70, dashed] (u) -- (w3);

    % Caixa do conjunto A_u = V(G) \ N(u)
    \draw[rounded corners=8pt, draw=blue!75!black, thick, fill=blue!4] (0, -0.8) rectangle (6.0, 0.4);
    \node[draw=none, fill=none, anchor=west] at (6.2, -0.2) {\small $A_u = V(G) \setminus N(u)$ ($|A_u| = n - d(u)$)};

    % Vértices em A_u
    \node[draw=blue!80!black, thick, fill=white] (v1) at (1.5, -0.2) {$v_1$};
    \node[draw=blue!80!black, thick, fill=white] (v2) at (3.2, -0.2) {$v_2$};
    \node[draw=blue!80!black, thick, fill=white] (v3) at (4.8, -0.2) {$v_3$};

    % Aresta interna em A_u
    \draw[very thick, blue!80!black] (v1) -- node[below, draw=none, fill=none, font=\scriptsize] {$e(A_u)$} (v2);

    % Arestas cruzadas e(A_u, B_u) (corte bipartido)
    \draw[thick, black!65] (w1) -- (v1);
    \draw[thick, black!65] (w2) -- (v1);
    \draw[thick, black!65] (w2) -- (v3);
    \draw[thick, black!65] (w3) -- (v3);
    \node[draw=none, fill=none, anchor=west, font=\scriptsize] at (6.2, 0.9) {Corte bipartido: $e(A_u, B_u) \le e(G) - t$};
\end{tikzpicture}
\end{center}

\subsubsection{Parte I --- Passo a Passo da Análise Local na Vizinhança de $u$:}
\begin{enumerate}
    \item \textbf{Contagem de Triângulos via Vizinhanças:} Seja $T(G)$ o número de triângulos em $G$. Toda aresta interna à vizinhança de um vértice $u$ forma um triângulo com $u$. Como cada triângulo possui três vértices, somar as arestas internas às vizinhanças de todos os vértices resulta em exatamente $3T(G)$.
    \item \textbf{Bipartição Natural centrada em $u$:} Para cada $u \in V(G)$, defina a partição $B_u = N(u)$ e $A_u = V(G) \setminus N(u)$, com cardinalidades $|B_u| = d(u)$ e $|A_u| = n - d(u)$. As arestas de $G$ dividem-se em três conjuntos disjuntos:
    \[ e(G) = e(A_u) + e(B_u) + e(A_u, B_u) \implies e(A_u) + e(B_u) = e(G) - e(A_u, B_u). \]
    \item \textbf{Hipótese de Distância à Bipartição ($t$-longe):} O subgrafo formado pelas arestas entre $A_u$ e $B_u$ é bipartido. Como $G$ é $t$-longe de ser bipartido, qualquer subgrafo bipartido possui no máximo $e(G) - t$ arestas, logo $e(A_u, B_u) \le e(G) - t \iff e(G) - e(A_u, B_u) \ge t$. Combinando com a igualdade anterior, obtemos:
    \[ e(A_u) + e(B_u) \ge t \implies e(A_u) \ge t - e(B_u). \]
    \item \textbf{Soma dos Graus em $A_u$:} Ao somar os graus (em $G$) dos vértices $v \in A_u$, cada aresta interna a $A_u$ contribui com $2$ e cada aresta entre $A_u$ e $B_u$ contribui com $1$. Substituindo $e(A_u, B_u) = e(G) - e(A_u) - e(B_u)$, temos:
    \[ \sum_{v \in A_u} d(v) = 2e(A_u) + e(A_u, B_u) = 2e(A_u) + e(G) - e(A_u) - e(B_u) = e(G) + e(A_u) - e(B_u). \]
    \item \textbf{Isolamento Local de $2e(B_u)$:} Aplicando a desigualdade $e(A_u) \ge t - e(B_u)$ na expressão da soma dos graus em $A_u$, eliminamos $e(A_u)$ e isolamos $2e(B_u)$:
    \[ \sum_{v \in A_u} d(v) \ge e(G) + \big(t - e(B_u)\big) - e(B_u) = e(G) + t - 2e(B_u) \implies 2e(B_u) \ge e(G) + t - \sum_{v \in A_u} d(v). \]
\end{enumerate}

\subsubsection{Parte II --- Passo a Passo da Agregação Global e Dupla Contagem:}
\begin{enumerate}[resume]
    \item \textbf{Somatório Global sobre $u \in V(G)$:} Somando a desigualdade local para todos os $n$ vértices $u \in V(G)$, estabelecemos a inequação global:
    \[ 2\sum_{u \in V(G)} e(B_u) \ge n\big(e(G) + t\big) - \sum_{u \in V(G)} \sum_{v \in A_u} d(v). \]
    \item \textbf{Inversão da Ordem do Somatório Duplo:} Em vez de fixar $u$ e somar $d(v)$ para $v \in A_u$, fixamos $v \in V(G)$ e contamos para quantos vértices $u$ vale a condição $v \in A_u$ (isto é, $v \notin N(u)$). Como existem exatamente $n - d(v)$ vértices que não são vizinhos de $v$, o somatório duplo reescreve-se como:
    \[ \sum_{u \in V(G)} \sum_{v \in A_u} d(v) = \sum_{v \in V(G)} d(v) \cdot |\{u \in V(G) : v \notin N(u)\}| = \sum_{v \in V(G)} d(v)\big(n - d(v)\big). \]
    \item \textbf{Cota via Desigualdade das Médias (MA-MG):} Pela desigualdade entre as médias aritmética e geométrica, temos $d(v)(n - d(v)) \le \frac{n^2}{4}$ para todo $v \in V(G)$. Somando sobre os $n$ vértices do grafo:
    \[ \sum_{u \in V(G)} \sum_{v \in A_u} d(v) = \sum_{v \in V(G)} d(v)\big(n - d(v)\big) \le \sum_{v \in V(G)} \frac{n^2}{4} = n \cdot \frac{n^2}{4}. \]
    \item \textbf{Substituição e Fechamento da Cota de Triângulos:} Substituindo o limite superior do somatório duplo na inequação global do Passo 6, obtemos:
    \[ 2\sum_{u \in V(G)} e(B_u) \ge n\big(e(G) + t\big) - n\frac{n^2}{4} \implies \sum_{u \in V(G)} e(B_u) \ge \frac{n}{2}\left(e(G) + t - \frac{n^2}{4}\right). \]
    Como a soma das arestas internas às vizinhanças equivale a $3T(G)$, dividimos ambos os lados por $3$ para concluir a prova:
    \[ T(G) = \frac{1}{3}\sum_{u \in V(G)} e(B_u) \ge \frac{n}{6}\left(e(G) + t - \frac{n^2}{4}\right). \]
\end{enumerate}

\textbf{Passo a Passo Rigoroso:}
\begin{enumerate}
    \item \textbf{Redução Local:} Para cada vértice $u$, separe o grafo em vizinhos ($B_u$) e não-vizinhos ($A_u$). O corte entre eles é bipartido, logo deixa pelo menos $t$ arestas no interior das partes: $e(A_u) + e(B_u) \ge t$.
    \item \textbf{Eliminação de $e(A_u)$:} Use a identidade $\sum_{v \in A_u} d(v) = 2e(A_u) + e(A_u, B_u) = e(G) + e(A_u) - e(B_u)$ junto com $e(A_u) \ge t - e(B_u)$ para isolar as arestas geradoras de triângulos: $2e(B_u) \ge e(G) + t - \sum_{v \in A_u} d(v)$.
    \item \textbf{Dupla Contagem e Maximização:} Some para todo $u \in V(G)$, inverta a soma dupla usando $\sum_{u}\sum_{v \in A_u} d(v) = \sum_{v} d(v)(n-d(v)) \le n\frac{n^2}{4}$ (por MA-MG) e divida por $6$ (fator $2$ de $2e(B_u)$ vezes fator $3$ de $3T(G)$).
\end{enumerate}

% EXERCÍCIO 7
\subsection*{Exercício 7: Teorema de Gerencsér-Gyárfás ($r(P_k) = \lceil 3k/2 \rceil$)}
\textbf{Intuição Estrutural:} O caminho $P_k$ possui comprimento $k$, logo contém $k+1$ vértices. A prova divide-se em duas frentes: (I) uma construção em blocos para o limite inferior que confina caminhos vermelhos à maior clique e força caminhos azuis a alternarem em um subgrafo bipartido; e (II) uma prova por contradição para o limite superior baseada no \textit{Lema de Cobertura de Caminhos}, onde maximizamos o caminho vermelho $R$ para forçar uma alta densidade de arestas azuis entre as extremidades do caminho azul $B$ e o interior de $R$.

\begin{center}
\begin{tikzpicture}[
    vertex/.style={circle, draw=black, fill=black, inner sep=2pt},
    blue edge/.style={blue, very thick},
    set label/.style={font=\Large\bfseries}
]

% Círculos vermelhos representando as arestas internas (os cliques monocromáticos)
\draw[red, ultra thick, fill=red!5] (0,0) circle (2.8cm);
\draw[red, ultra thick, fill=red!5] (7,0) circle (2cm);

% Vértices do Conjunto A (alinhados na borda direita para facilitar a visão do zigue-zague)
\node[vertex, label=left:{$a_1$}] (a1) at (1.5, 1.5) {};
\node[vertex, label=left:{$a_2$}] (a2) at (1.5, 0.5) {};
\node[vertex, label=left:{$a_3$}] (a3) at (1.5, -0.5) {};
\node[vertex, label=left:{$a_4$}] (a4) at (1.5, -1.5) {};
% Alguns vértices soltos para preencher visualmente o conjunto A
\node[vertex] (a5) at (-1, 0) {}; 
\node[vertex] (a6) at (0, 1.2) {};
\node[vertex] (a7) at (-0.5, -1.5) {};
\node[vertex] (a8) at (0, -0.2) {};

% Vértices do Conjunto B (alinhados na borda esquerda)
\node[vertex, label=right:{$b_1$}] (b1) at (5.5, 1) {};
\node[vertex, label=right:{$b_2$}] (b2) at (5.5, 0) {};
\node[vertex, label=right:{$b_3$}] (b3) at (5.5, -1) {};

% Caminho Azul intercalando A e B (Zigue-zague)
\draw[blue edge] (a1) -- (b1);
\draw[blue edge] (b1) -- (a2);
\draw[blue edge] (a2) -- (b2);
\draw[blue edge] (b2) -- (a3);
\draw[blue edge] (a3) -- (b3);
\draw[blue edge] (b3) -- (a4);

% Rótulos dos Conjuntos
\node[set label, above] at (0, 3.2) {Conjunto $A$};
\node[above] at (0, 2.8) {$|A| = k$};
\node[red, font=\bfseries, align=center] at (0, -3.5) {};

\node[set label, above] at (7, 2.5) {Conjunto $B$};
\node[above] at (7, 1.9) {$|B| = \lfloor (k-1)/2 \rfloor$};
\node[red, font=\bfseries, align=center] at (7, -2.7) {};

% Rótulo explicativo do caminho azul
\node[blue, font=\bfseries, fill=white, inner sep=4pt] at (3.5, 0) {};
\node[blue, fill=white, inner sep=2pt] at (3.5, -0.5) {};

\end{tikzpicture}
\end{center}

\subsubsection{Parte I --- Passo a Passo do Limite Inferior ($r(P_k) \ge \lceil 3k/2 \rceil$):}
\begin{enumerate}
    \item \textbf{Partição Extremal ($N = \lceil 3k/2 \rceil - 1$):} Divida os $N$ vértices de $K_N$ em dois conjuntos disjuntos $A$ e $B$ tais que $|A| = k$ e $|B| = \lfloor (k-1)/2 \rfloor$. A soma resulta exatamente em $k + \lfloor (k-1)/2 \rfloor = \lceil 3k/2 \rceil - 1 = N$.
    \item \textbf{Coloração em Blocos:} Pinte todas as arestas internas a $A$ e internas a $B$ de vermelho (duas cliques vermelhas disjuntas) e todas as arestas entre $A$ e $B$ de azul (um grafo bipartido completo azul).
    \item \textbf{Confinamento Vermelho:} Sem arestas vermelhas cruzando entre $A$ e $B$, qualquer caminho vermelho vive inteiramente em $A$ ou inteiramente em $B$. Como $|A| > |B|$, o maior caminho vermelho tem no máximo $|A| = k$ vértices.
    \item \textbf{Alternância Azul:} Como o subgrafo azul é bipartido entre $A$ e $B$, todo caminho azul alterna estritamente entre os dois conjuntos. Sendo $B$ o menor conjunto, o caminho azul visita no máximo os $|B|$ vértices de $B$ e no máximo $|B| + 1$ vértices de $A$, totalizando no máximo $2|B| + 1 = 2\lfloor (k-1)/2 \rfloor + 1 \le k$ vértices. Logo, nenhum caminho monocromático atinge $k+1$ vértices.
\end{enumerate}

\newpage

\textbf{O Lema de Cobertura de Caminhos (Gerencsér-Gyárfás):}

Antes de atacar o limite superior, precisamos deste lema crucial: \textit{"Os vértices de qualquer grafo completo $K_N$ com arestas bicoloridas podem ser particionados em exatamente dois caminhos disjuntos de cores distintas."}
\begin{quote}
\small
\textbf{Prova do Lema:} Seja $R = (r_1, \dots, r_a)$ um caminho vermelho e $B = (b_1, \dots, b_b)$ um caminho azul, disjuntos, que juntos cubram o \textit{maior número possível} de vértices. Suponha, por absurdo, que a cobertura não seja total e exista um vértice $x \notin R \cup B$.

A aresta $x r_1$ não pode ser vermelha (caso contrário, poderíamos estender $R$ para $x \to r_1 \dots$, cobrindo $x$ e contradizendo a maximalidade da nossa escolha inicial). Logo, $x r_1$ é azul. Analogamente, a aresta $x b_1$ não pode ser azul (caso contrário, estenderíamos $B$ para $x \to b_1 \dots$), o que obriga $x b_1$ a ser vermelha.

Agora, avalie a aresta conectora $r_1 b_1$ entre os caminhos originais:
\begin{itemize}
    \item \textbf{Se $r_1 b_1$ for vermelha:} O caminho $x \to u_1 \to v_1 \to v_2 \dots \to v_r$ é um novo caminho vermelho longo. O caminho azul restante é $B' = (b_2, \dots, b_b)$. Essa nova configuração é disjunta e cobre $a+b+1$ vértices, pois absorveu o vértice esquecido $x$, violando a suposição de maximalidade inicial.
    \item \textbf{Se $r_1 b_1$ for azul:} O caminho $x \to v_1 \to u_1 \to u_2 \dots \to u_b$ é um novo caminho azul longo. O caminho vermelho restante é $R' = (v_2, \dots, v_r)$. Assim como no caso anterior, cobrimos todos os vértices e mais o $x$, gerando absurdo!
\end{itemize}
Como ambos os casos levam a contradições insolúveis, a hipótese de haver um vértice $x$ de fora é falsa. Fica provado que podemos sempre particionar o conjunto de vértices integralmente em exatamente dois caminhos de cores opostas. $\qed$
\end{quote}
\begin{center}
    \begin{tikzpicture}[
        % Definição dos estilos visuais para os nós e arestas
        vertex/.style={circle, draw=black, thick, fill=white, inner sep=2pt, minimum size=8pt},
        red edge/.style={red, line width=1.8pt},
        blue edge/.style={blue, line width=1.8pt},
        red dashed edge/.style={red, thick, dashed, line width=1.2pt},
        blue dashed edge/.style={blue, thick, dashed, line width=1.2pt},
        query edge/.style={black, thick, dotted, line width=1.2pt}
    ]

    % Nodos do caminho vermelho (R) na parte superior
    \node[vertex, label=above:{\Large $v_1$}] (v1) at (3, 2) {};
    \node[vertex, label=above:{\Large $v_2$}] (v2) at (5, 2) {};
    \node at (7, 2) {\Large $\dots$}; % Elipse indicando continuação
    \node[vertex, label=above:{\Large $v_r$}] (vr) at (9, 2) {};
    
    % Nodos do caminho azul (B) na parte inferior
    \node[vertex, label=below:{\Large $u_1$}] (u1) at (3, -1) {};
    \node[vertex, label=below:{\Large $u_2$}] (u2) at (5, -1) {};
    \node at (7, -1) {\Large $\dots$}; % Elipse indicando continuação
    \node[vertex, label=below:{\Large $u_b$}] (ub) at (9, -1) {};
    
    % Nodo isolado x, centralizado verticalmente à esquerda
    \node[vertex, label=left:{\Large $x$}] (x) at (1, 0.5) {};
    
    % Arestas do caminho vermelho
    \draw[red edge] (v1) -- (v2);
    \draw[red edge] (v2) -- (6.5, 2);
    \draw[red edge] (7.5, 2) -- (vr);
    
    % Arestas do caminho azul
    \draw[blue edge] (u1) -- (u2);
    \draw[blue edge] (u2) -- (6.5, -1);
    \draw[blue edge] (7.5, -1) -- (ub);
    
    % Arestas tracejadas indicando as conexões forçadas por maximalidade
    \draw[blue dashed edge] (x) -- (v1);
    \draw[red dashed edge] (x) -- (u1);
    
    % Aresta pontilhada central (o ponto da contradição)
    \draw[query edge] (v1) -- node[right=2pt] {\Large ?} (u1);
    
    % Rótulos de texto dos caminhos
    \node[above] at (6, 2.6) {\large $R$ (Caminho Vermelho)};
    \node[below] at (6, -1.6) {\large $B$ (Caminho Azul)};
    
    \end{tikzpicture}
\end{center}


\newpage

\subsubsection{Parte II --- Passo a Passo do Limite Superior ($r(P_k) \le \lceil 3k/2 \rceil$):}
\begin{enumerate}[resume]
    \item \textbf{Hipótese por Absurdo e Lema de Cobertura:} Tome $N = \lceil 3k/2 \rceil$ e suponha que exista uma 2-coloração de $K_N$ sem $P_k$ monocromático (sem caminhos de $k+1$ vértices). Pelo Lema de Gerencsér-Gyárfás, $V(K_N)$ pode ser particionado em um caminho vermelho $R = (v_1, \dots, v_r)$ e um caminho azul $B = (u_1, \dots, u_b)$ disjuntos, com $r + b = N$.
    \item \textbf{Escolha Maximal e Cota para $B$:} Escolha a partição que \textit{maximiza} o tamanho $r$ do caminho vermelho $R$. Pela hipótese de absurdo, $r \le k$ e $b \le k$. Disso decorre o limite inferior para o tamanho de $B$:
    \[ b = N - r \ge \left\lceil \frac{3k}{2} \right\rceil - k = \left\lceil \frac{k}{2} \right\rceil. \]
    \item \textbf{Bloqueio nas Pontas ($v_r u_1$ é Azul):} A aresta $v_r u_1$ não pode ser vermelha; se fosse, estenderíamos $R$ para $(v_1, \dots, v_r, u_1)$ com tamanho $r+1$ e reduziríamos $B$ para $(u_2, \dots, u_b)$, violando a maximalidade de $r$. Logo, $v_r u_1$ é obrigatoriamente azul.
    \item \textbf{Bloqueio de Inserção no Interior de $R$:} Não é possível inserir $u_1$ entre dois vértices adjacentes $v_i$ e $v_{i+1}$ de $R$ usando duas arestas vermelhas, pois isso também aumentaria $R$ para $r+1$ vértices. Assim, para cada par consecutivo $(v_i, v_{i+1})$ em $R$, pelo menos uma das arestas ligadas a $u_1$ é azul.
    \item \textbf{Contagem de Arestas Azuis nas Extremidades de $B$:} Combinando os dois passos anteriores, o vértice $u_1$ incide em pelo menos $\lceil r/2 \rceil$ arestas azuis conectadas a $R$. Simetricamente, o vértice final $u_b$ também incide em pelo menos $\lceil r/2 \rceil$ arestas azuis ligadas a $R$.
    \item \textbf{Fusão por Atalhos Azuis e Contradição:} Com as extremidades de $B$ maciçamente conectadas a $R$ por arestas azuis, forma-se uma estrutura cíclica que permite mesclar todos os $b$ vértices de $B$ com pelo menos $\lfloor r/2 \rfloor + 1$ vértices de $R$ através de atalhos azuis. Usando a identidade $r - \lfloor r/2 \rfloor = \lceil r/2 \rceil$, o tamanho desse novo caminho azul é:
    \[ b + \left\lfloor \frac{r}{2} \right\rfloor + 1 = N - r + \left\lfloor \frac{r}{2} \right\rfloor + 1 = N - \left\lceil \frac{r}{2} \right\rceil + 1. \]
    Substituindo $N = k + \lceil k/2 \rceil$ e usando que $r \le k \implies \lceil r/2 \rceil \le \lceil k/2 \rceil$, obtemos no mínimo:
    \[ \left(k + \left\lceil \frac{k}{2} \right\rceil\right) - \left\lceil \frac{k}{2} \right\rceil + 1 = k + 1 \text{ vértices.} \]
    Isso força um caminho azul de tamanho $k+1$, contradizendo a premissa inicial e provando que $r(P_k) = \lceil 3k/2 \rceil$.
\end{enumerate}

% EXERCÍCIOS 8 E 9
\subsection*{Exercício 8 e 9: Independência de Variáveis Aleatórias}
\textbf{Ex 8 (Contraexemplo):} Independência dois a dois não garante independência mútua. O clássico lançamento de duas moedas com eventos $A$ (moeda 1 coroa), $B$ (moeda 2 coroa), $C$ (mesmo resultado). $P(A \cap B \cap C) = 1/4 \ne P(A)P(B)P(C) = 1/8$.

\textbf{Ex 9 (Produto de VA Independentes):} Demonstração direta via Fubini e fatoração da densidade conjunta. A independência quebra a integral dupla $\iint xy f(x,y) dx dy$ em duas integrais marginais cujo produto é exatamente $\mathbb{E}[X_1]\mathbb{E}[X_2]$.

\newpage

% EXERCÍCIO 10
\subsection*{Exercício 10: Limite do Número de Independência por Probabilidade}
\textbf{Intuição Probabilística:} Como achar um conjunto independente grande sem construir heurísticas complexas? Sorteie a ordem dos vértices e retenha os que aparecem "primeiro" entre todos os seus vizinhos. Esse grupo é, por definição, independente.

\begin{center}
\begin{tikzpicture}[scale=0.9, every node/.style={circle, draw, minimum size=6mm, inner sep=1pt}]
    \node[fill=gray!30] (a) at (0,0) {$v_a$};
    \node[fill=gray!30] (b) at (2,0) {$v_b$};
    \node[fill=gray!30] (c) at (4,0) {$v_c$};
    \node (d) at (6,0) {$v_d$};
    \node (e) at (8,0) {$v_e$};

    \draw[thick, ->] (a) edge[bend left=30] (d);
    \draw[thick, ->] (b) edge[bend left=30] (e);
    \draw[thick, ->] (c) edge[bend left=30] (d);
    
    \node[draw=none, yshift=-0.7cm] at (a) {$\sigma=1$};
    \node[draw=none, yshift=-0.7cm] at (d) {$\sigma=4$};
\end{tikzpicture}
\end{center}

\textbf{Passo a Passo Rigoroso:}
\begin{enumerate}
    \item Defina uma ordem aleatória uniforme $\sigma$ de $V$. O evento "v é o menor na sua vizinhança" garante independência. (Uma aresta entre os guardados implicaria que um deles tem um vizinho menor).
    \item O conjunto resultante $A$ tem esperança igual à soma das probabilidades indicadoras.
    \item $P(v \in A) = \frac{1}{d(v)+1}$, pois todos os elementos da vizinhança fechada $N[v]$ têm chance idêntica de cair na primeira posição.
    \item Como $\alpha(G) \ge |A|$ em \textit{qualquer} realização, ele supera a média: $\alpha(G) \ge \mathbb{E}[|A|] = \sum \frac{1}{d(v)+1}$.
\end{enumerate}

% EXERCÍCIO 11
\subsection*{Exercício 11: Desigualdade de Jensen}
\textbf{Intuição Geométrica e Probabilística:} A Desigualdade de Jensen formaliza que a média de uma função convexa é maior ou igual à função aplicada na média. Geometricamente, uma função convexa "faz uma curva para cima", o que significa que qualquer reta tangente traçada em um ponto de seu gráfico ficará inteiramente abaixo da função. Ao escolhermos o ponto de tangência exatamente no valor esperado da variável $\mathbb{E}[X]$, a linearidade da esperança anula a inclinação da reta, provando o resultado.

\textbf{Passo a Passo da Demonstração (Desigualdade de Jensen):}
\begin{enumerate}
    \item \textbf{Ponto de Tangência e Aproximação de Taylor:} Seja $x_0 = \mathbb{E}[X]$ o ponto onde avaliaremos a função. Assumindo (para simplificação didática) que $f$ é duas vezes diferenciável, a fórmula de Taylor com resto de Lagrange garante que para qualquer $x \in \mathbb{R}$, existe um $\xi$ entre $x$ e $x_0$ tal que:
    \[ f(x) = f(x_0) + f'(x_0)(x - x_0) + f''(\xi)\frac{(x-x_0)^2}{2}. \]
    \item \textbf{Limitante Inferior via Convexidade:} Como a função $f$ é convexa, sua segunda derivada é não negativa ($f''(\xi) \ge 0$). Isso implica que o termo de erro na expansão de Taylor é positivo ou nulo. Ao removê-lo, obtemos a equação da reta tangente $L(x) := f(x_0) + f'(x_0)(x - x_0)$, que atua como um limite inferior global:
    \[ f(x) \ge L(x) \quad \text{para todo } x \in \mathbb{R}. \]
    \textit{Nota:} Se $f$ não for duas vezes diferenciável, a convexidade por si só garante a existência de uma reta suporte (subderivada) que cumpre exatamente essa mesma propriedade.
    \item \textbf{Aplicação da Variável Aleatória e Monotonicidade:} Substituindo a variável real $x$ pela variável aleatória $X$, a desigualdade se mantém evento a evento: $f(X) \ge L(X)$. Aplicando a esperança matemática em ambos os lados (uma operação que preserva desigualdades), obtemos:
    \[ \mathbb{E}[f(X)] \ge \mathbb{E}[L(X)]. \]
    \item \textbf{Linearidade e Cancelamento Mágico:} Como a função $L(X)$ é afim (linear), aplicamos a linearidade da esperança no lado direito:
    \[ \mathbb{E}[L(X)] = f(x_0) + f'(x_0)\big(\mathbb{E}[X] - x_0\big). \]
    Pela nossa escolha inicial no Passo 1, temos $x_0 = \mathbb{E}[X]$. Logo, a diferença $\big(\mathbb{E}[X] - x_0\big)$ é exatamente zero, o que anula o termo da derivada.
    \item \textbf{Conclusão:} O cálculo do lado direito colapsa simplesmente para $f(x_0)$. Substituindo $x_0$ de volta por $\mathbb{E}[X]$, concluímos a prova:
    \[ \mathbb{E}[f(X)] \ge f(\mathbb{E}[X]). \]
\end{enumerate}

% EXERCÍCIO 12
\subsection*{Exercício 12: Supersaturação}
\textbf{Intuição Estrutural e Algébrica:} A supersaturação fraca garante que, se ultrapassarmos significativamente o limite extremal de um subgrafo $H$ (neste caso, o dobro de arestas), não apenas uma, mas \textit{muitas} cópias de $H$ devem aparecer inevitavelmente. A prova usa um argumento construtivo de remoção por absurdo: se houvessem "poucas" cópias, poderíamos removê-las tirando pouquíssimas arestas, resultando em um grafo $H$-livre que ainda teria "arestas demais" (violando a própria definição fundamental do número extremal).

\textbf{Passo a Passo da Demonstração (Supersaturação Fraca):}
\begin{enumerate}
    \item \textbf{Hipótese por Absurdo:} Considere um grafo $G$ com $n$ vértices e número de arestas $e(G) \ge 2 \cdot ex(n, H)$. Suponha, por contradição, que $G$ contenha um número $k$ de cópias de $H$ estritamente menor que $ex(n, H)$, isto é, $k < ex(n, H)$.
    \item \textbf{Destruição das Cópias (Remoção Estratégica):} Para transformar $G$ em um grafo totalmente livre de $H$, basta remover ao menos uma aresta de cada uma das $k$ cópias existentes. O número máximo de arestas que precisaremos remover para limpar o grafo é exatamente $k$ (o pior cenário, onde as cópias são totalmente disjuntas em relação às suas arestas).
    \item \textbf{Construção e Limite Inferior de $G'$:} Seja $G'$ o grafo resultante após remover essas $k$ (ou menos) arestas. O número de arestas restantes em $G'$ obedece necessariamente ao seguinte limite inferior: 
    \[ e(G') \ge e(G) - k. \]
    \item \textbf{Definição Extremal (Limite Superior):} Por construção, $G'$ é um grafo $H$-livre com $n$ vértices. Pela definição clássica de número extremal, qualquer grafo com essas propriedades pode ter, no máximo, $ex(n, H)$ arestas. Logo, aplica-se o limite superior ao novo grafo: 
    \[ e(G') \le ex(n, H). \]
    \item \textbf{Contradição Algébrica:} Juntando o limite inferior e o limite superior de $e(G')$, deduzimos que $e(G) - k \le ex(n, H)$. Substituindo a premissa inicial $e(G) \ge 2 \cdot ex(n, H)$ na inequação, obtemos:
    \[ 2 \cdot ex(n, H) - k \le ex(n, H). \]
    Isolando a variável $k$ algebricamente, chegamos a $ex(n, H) \le k$, o que contradiz diretamente a nossa suposição inicial de que $k < ex(n, H)$. Como a suposição leva a um absurdo lógico, o grafo $G$ deve conter pelo menos $ex(n, H)$ cópias de $H$.
\end{enumerate}

% EXERCÍCIO 13
\subsection*{Exercício 13: Teorema de Turán via Probabilidade}
\textbf{Intuição Estrutural e Algébrica:} Conecta-se grafos complementares, o Exercício 10 e a desigualdade de Cauchy-Schwarz.

\textbf{Passo a Passo Rigoroso:}
\begin{enumerate}
    \item $G$ é livre de $K_r \implies$ o complementar $\overline{G}$ tem número de independência $< r$.
    \item Pelo Ex 10, $\alpha(\overline{G}) \ge \sum \frac{1}{1+\overline{d}(v)} = \sum \frac{1}{n-d(v)}$. Portanto $r-1 \ge \sum \frac{1}{n-d(v)}$.
    \item Usa-se Cauchy-Schwarz no somatório com os termos $(n-d(v))$. O somatório cruzado gera $(n^2 - 2e(G))$. Isola-se o termo de arestas, recuperando exatamente o limitante clássico da densidade de Turán.
\end{enumerate}

% EXERCÍCIO 14
\subsection{Exercício 14: Grafos Aleatórios (Monotonicidade e Acoplamento)}
\textbf{Intuição Probabilística:} Para provar que a probabilidade de $G(n,p)$ possuir uma propriedade crescente aumenta junto com $p$, utilizamos a técnica de \textit{Acoplamento (Coupling)}. Em vez de gerar os grafos $G(n,p)$ e $G(n,p')$ em universos isolados, geramos ambos simultaneamente no mesmo espaço de probabilidade usando variáveis uniformes. Isso cria uma dependência estrutural direta onde o grafo com menor probabilidade torna-se um subgrafo explícito do grafo com maior probabilidade.

\textbf{Passo a Passo Rigoroso:}

\subsubsection{Parte I --- Construção do Espaço de Probabilidade Acoplado:}
\begin{enumerate}
    \item \textbf{Definição do Objetivo:} Seja $\mathcal{A}$ uma família crescente de grafos. Queremos mostrar que para quaisquer probabilidades $p, p'$ tais que $0 \le p \le p' \le 1$, a função $f(p) = \mathbb{P}(G(n,p) \in \mathcal{A})$ obedece $f(p) \le f(p')$.
    \item \textbf{Variáveis Uniformes Independentes:} Considere o grafo completo $K_n$, que possui $N = \binom{n}{2}$ arestas. Para cada aresta $e \in E(K_n)$, associe uma variável aleatória independente $U_e \sim \text{Uniforme}[0, 1]$.
    \item \textbf{Geração Simulada de $G_p$:} Para um dado limiar $p \in [0,1]$, defina o grafo aleatório $G_p$ selecionando apenas as arestas cujo valor sorteado não excede $p$: 
    \[ G_p := \{e \in E(K_n) : U_e \le p\}. \]
    \item \textbf{Validação da Distribuição:} Como a variável $U_e$ é uniforme no intervalo $[0,1]$, a probabilidade de uma aresta existir em $G_p$ é exatamente $\mathbb{P}(U_e \le p) = p$. Além disso, as arestas são independentes. Logo, $G_p$ possui exatamente a mesma distribuição do modelo de Erdős-Rényi $G(n,p)$, o que implica que $f(p) = \mathbb{P}(G_p \in \mathcal{A})$.
\end{enumerate}

\subsubsection{Parte II --- Monotonicidade Estrutural e Inclusão de Eventos:}
\begin{enumerate}[resume]
    \item \textbf{Inclusão de Subgrafos ($G_p \subseteq G_{p'}$):} Considere agora $p \le p'$. Para qualquer realização das variáveis aleatórias $U_e$, se uma aresta $e$ pertence a $G_p$, então $U_e \le p$. Como $p \le p'$, transitivamente temos $U_e \le p'$, o que significa que a aresta $e$ obrigatoriamente também pertence a $G_{p'}$. Portanto, $G_p \subseteq G_{p'}$ para todo evento do espaço amostral.
    \item \textbf{Aplicação da Propriedade Crescente:} A família $\mathcal{A}$ é, por definição, fechada para a adição de arestas. Isso significa que, se um grafo possui a propriedade e adicionamos arestas a ele, o novo grafo continua possuindo a propriedade. Logo, como $G_p \subseteq G_{p'}$, a implicação lógica $G_p \in \mathcal{A} \implies G_{p'} \in \mathcal{A}$ é sempre verdadeira.
    \item \textbf{Inclusão de Eventos:} A implicação determinística acima traduz-se na linguagem de Teoria da Medida como uma inclusão de conjuntos (eventos) no espaço amostral: 
    \[ \{G_p \in \mathcal{A}\} \subseteq \{G_{p'} \in \mathcal{A}\}. \]
    \item \textbf{Cálculo da Probabilidade e Conclusão:} Pela propriedade fundamental da monotonicidade das medidas de probabilidade (se $X \subseteq Y$, então $\mathbb{P}(X) \le \mathbb{P}(Y)$), aplicamos a probabilidade em ambos os lados da inclusão de eventos:
    \[ f(p) = \mathbb{P}(G_p \in \mathcal{A}) \le \mathbb{P}(G_{p'} \in \mathcal{A}) = f(p'). \]
    Como $0 \le p \le p' \le 1 \implies f(p) \le f(p')$, concluímos que a função $f(p)$ é monótona não decrescente (crescente no sentido amplo).
\end{enumerate}



\subsection*{Exercício 15: Grafos Aleatórios (Threshold de Grau Minímo)}

\textbf{Intuição Probabilística:} O problema analisa a evolução do grau mínimo de um grafo aleatório $G(n,p)$. A variável chave é $X$, o número de vértices com grau "proibido" (menor que $d$). A demonstração explora a transição de fase dessa variável: se a esperança $\mathbb{E}[X]$ tende a zero, usamos o Método do Primeiro Momento (Markov) para provar que $X=0$ com alta probabilidade; se a esperança tende ao infinito, usamos o Método do Segundo Momento (Chebyshev e variância) para provar que $X > 0$ com alta probabilidade.

\textbf{Passo a Passo Rigoroso:}
\begin{enumerate}
    \item \textbf{O Valor Esperado ($\mathbb{E}[X_k]$) via Variáveis Indicadoras:} 
    A variável $X_k$ representa o número total de vértices no grafo que possuem grau exatamente igual a $k$. Para calcular sua esperança, não buscamos o grau médio de um vértice (que seria $\mathbb{E}[d_G(v)] = (n-1)p$), mas sim "quantos vértices esperamos encontrar com grau $k$".
    \begin{itemize}
        \item \textit{Probabilidade do evento exato:} Para que um vértice $v$ tenha grau $k$, ele precisa de exatamente $k$ sucessos em $n-1$ tentativas (arestas possíveis). A chance disso ocorrer é dada pela massa da distribuição binomial:
        \[ \Pr(d_G(v) = k) = \binom{n-1}{k} p^k (1-p)^{n-1-k} \]
        \item \textit{Linearidade da Esperança:} Definimos uma variável indicadora $I_v$ que vale $1$ se o vértice $v$ tem grau $k$, e $0$ caso contrário. O total de vértices com grau $k$ é a soma desses indicadores: $X_k = \sum_{v \in V} I_v$.
        \item \textit{Soma global:} Como existem $n$ vértices com probabilidades idênticas, a soma dos valores esperados individuais é:
        \[ \mathbb{E}[X_k] = \sum_{v \in V} \mathbb{E}[I_v] = \sum_{v \in V} \Pr(d_G(v) = k) = n \binom{n-1}{k} p^k (1-p)^{n-1-k} \]
    \end{itemize}
    
    \item \textbf{Aproximação Assintótica e Dominação do Termo $X_{d-1}$:} 
     Para lidar com o comportamento do grafo quando $n \to \infty$, analisamos e simplificamos os componentes de $\mathbb{E}[X_k]$ separadamente:
    \begin{itemize}
        \item limiar crítico fornecido pelo teorema é $p = \frac{\log n + (d-1)\log\log n \pm \omega(n)}{n}$. Como o denominador $n$ cresce de forma muito mais acelerada que o numerador (que é dominado pelo termo $\log n$), a probabilidade de uma aresta existir tende a zero quando $n \to \infty$. Matematicamente, denotamos essa tendência por $p = o(1)$.
        
        \item Sabendo que $p \to 0$ e que $k$ é uma constante, o termo combinatório simplifica-se para $\binom{n-1}{k} \sim \frac{n^k}{k!}$. Além disso, usando o limite notável da função exponencial, o termo de probabilidade para a não-existência de arestas aproxima-se por $(1-p)^{n-1-k} \sim e^{-np}$.
        
        \item O termo $e^{-np}$ é o principal responsável pelo decaimento probabilístico. Substituindo a expressão de $p$ diretamente no expoente, o fator $n$ se cancela:
        \[ e^{-np} \sim \exp\big(-(\log n + (d-1)\log\log n \pm \omega(n))\big). \]
        Usando as propriedades básicas de exponenciais e logaritmos ($e^{-a-b} = e^{-a}e^{-b}$ e $e^{-\log x} = \frac{1}{x}$), essa expressão colapsa nitidamente para:
        \[ e^{-np} \sim \frac{1}{n} \cdot \frac{1}{(\log n)^{d-1}} \cdot e^{\mp\omega(n)}. \]
        
        \item Juntando a expansão acima com as aproximações iniciais (e notando que $p^k \sim \left(\frac{\log n}{n}\right)^k$), a esperança de encontrarmos um vértice de grau $k$ assume a forma \[\mathbb{E}[X_k] \sim \frac{(\log n)^{k-d+1}}{k!}e^{\mp\omega(n)} \quad .\] Observe o expoente do logaritmo ($k-d+1$): para qualquer grau $k < d-1$, esse expoente será negativo, forçando $\mathbb{E}[X_k]$ a tender rapidamente a zero. Portanto, a esperança total de vértices com grau "proibido" (grau $< d$), $\mathbb{E}[X] = \sum_{k=0}^{d-1} \mathbb{E}[X_k]$, é assintoticamente dominada pelo seu maior termo, $k = d-1$, onde o expoente logarítmico desaparece (fica elevado a zero). Isso nos entrega a forma final:
        \[ \mathbb{E}[X] \sim \mathbb{E}[X_{d-1}] \sim \frac{1}{(d-1)!}e^{\mp\omega(n)}. \]
    \end{itemize}


    \item \textbf{Método do Primeiro Momento ($p$ acima do limiar):} 
     Para um $p$ grande (sinal $+\omega(n)$), temos $e^{-\omega(n)} \to 0$, resultando em $\mathbb{E}[X] \to 0$. Para transformar essa esperança nula em uma garantia probabilística, usamos a Desigualdade de Markov:
    \begin{itemize}
        \item Para qualquer variável aleatória não negativa $X$ e limite $a > 0$, Markov garante que $\mathbb{P}(X \ge a) \le \frac{\mathbb{E}[X]}{a}$. Nossa variável $X$ conta vértices, logo é não negativa. Queremos medir a chance de existir \textit{pelo menos um} vértice com grau proibido, então aplicamos $a = 1$.
        \item \textit{O Limite Superior:} Substituindo na fórmula, a probabilidade do grafo falhar fica limitada pela sua própria média:
        \[ \mathbb{P}(X \ge 1) \le \frac{\mathbb{E}[X]}{1} = \mathbb{E}[X]. \]
        Como demonstramos que $\mathbb{E}[X] \to 0$, concluímos inevitavelmente que $\mathbb{P}(X \ge 1) \to 0$.
        \item  O evento $\{X \ge 1\}$ significa "existe ao menos um vértice com grau menor que $d$". Seu evento complementar é $\{X = 0\}$, que significa "nenhum vértice tem grau menor que $d$". Se nenhum vértice tem grau menor que $d$, então todos os vértices têm grau maior ou igual a $d$. Essa é a definição exata de o grau mínimo do grafo ser $\delta(G) \ge d$.
        \item A probabilidade do evento complementar é $1$ subtraído da probabilidade do evento original. Logo, o grau mínimo atinge o limiar esperado com probabilidade $1$:
        \[ \mathbb{P}(\delta(G) \ge d) = \mathbb{P}(X = 0) = 1 - \mathbb{P}(X \ge 1) \to 1 - 0 = 1. \]
    \end{itemize}

    \item \textbf{Método do Segundo Momento ($p$ abaixo do limiar):} 
   Para um $p$ pequeno (sinal $-\omega(n)$), a esperança diverge: $\mathbb{E}[X] \to \infty$. Para provar que o grafo realmente possui vértices com grau proibido ($X > 0$), usamos o Método do Segundo Momento via Desigualdade de Chebyshev:
    \begin{itemize}
        \item \textit{A Desigualdade de Chebyshev:} Esta desigualdade limita a probabilidade de uma variável se afastar de sua média ($\mathbb{P}(|X - \mathbb{E}[X]| \ge t) \le \frac{\text{Var}(X)}{t^2}$). Queremos a probabilidade de o grafo falhar, ou seja, não ter nenhum vértice de grau menor que $d$ ($X = 0$). Se $X=0$, o afastamento em relação à média é exatamente $t = \mathbb{E}[X]$. Substituindo na fórmula, obtemos o limite clássico:
        \[ \mathbb{P}(X = 0) \le \frac{\text{Var}(X)}{\mathbb{E}[X]^2}. \]
        \item \textit{Transformando a Variância:} Pela definição matemática de variância, sabemos que $\text{Var}(X) = \mathbb{E}[X^2] - \mathbb{E}[X]^2$. Dividindo os dois termos por $\mathbb{E}[X]^2$, a nossa desigualdade se reescreve como:
        \[ \mathbb{P}(X=0) \le \frac{\mathbb{E}[X^2]}{\mathbb{E}[X]^2} - 1. \]
         O termo $\mathbb{E}[I_u I_v]$ é a probabilidade conjunta de os vértices $u$ e $v$ falharem simultaneamente. O grau de $u$ depende das suas $n-1$ arestas possíveis, e o de $v$ das dele. A \textit{única} aresta do universo que eles compartilham e que poderia "vincular" o destino dos dois é a aresta $\{u,v\}$. Num grafo gigantesco ($n \to \infty$), o impacto de uma única aresta é irrisório. Assim, os eventos comportam-se de forma quase independente: $\mathbb{E}[I_u I_v] \sim \mathbb{E}[I_u]\mathbb{E}[I_v]$.
        \item \textit{A Razão Final:} Assumindo a independência, a soma dos pares cruzados se aproxima de $\mathbb{E}[X]^2$. Isso nos dá a relação direta:
        \[ \mathbb{E}[X^2] \sim \mathbb{E}[X] + \mathbb{E}[X]^2. \]
        Por fim, dividimos toda essa equação por $\mathbb{E}[X]^2$. O lado direito se divide em duas frações simples, revelando a aproximação procurada:
        \[ \frac{\mathbb{E}[X^2]}{\mathbb{E}[X]^2} \sim \frac{\mathbb{E}[X]}{\mathbb{E}[X]^2} + \frac{\mathbb{E}[X]^2}{\mathbb{E}[X]^2} = \frac{1}{\mathbb{E}[X]} + 1. \]
        \item \textit{Cancelamento Algébrico:} Substituindo essa convergência na nossa desigualdade de probabilidade, o termo $1$ se cancela perfeitamente com o $-1$ vindo da variância:
        \[ \mathbb{P}(X=0) \le \left(\frac{1}{\mathbb{E}[X]} + 1\right) - 1 = \frac{1}{\mathbb{E}[X]}. \]
        \item \textit{Conclusão do Limiar:} Como calculamos anteriormente que $\mathbb{E}[X] \to \infty$, a fração $\frac{1}{\mathbb{E}[X]}$ inevitavelmente tende a $0$. Assim, $\mathbb{P}(X = 0) \to 0$. Como o evento $\{X = 0\}$ equivale a dizer que o grau mínimo de todos os vértices é $\ge d$, fica provado que, abaixo do limiar, $\mathbb{P}(\delta(G) \ge d) \to 0$.
    \end{itemize}
\end{enumerate}


\end{document}